\section{Discussion}

\label{sec:discussion}

\subsection{Parameter Selection}

The choice of $\lambda$, $\mu_i$, and $\bm{w}$ significantly affects market behavior. We recommend:
\begin{itemize}[leftmargin=1.1em]
    \item $\lambda \in [0.5, 5.0]$: Lower for volatile resource pairs, higher for stable pairs.
    \item $\mu_i \in [0.01, 0.1]$: Provides inertia without excessive price lag.
    \item $w_i$ proportional to the economic value of resource $i$ in the target market.
\end{itemize}

Governance can adjust these parameters via time-locked proposals with a 7-day voting period.

\subsection{Limitations}

\begin{enumerate}[leftmargin=1.4em]
    \item \textbf{Computational overhead}: Ruth-4 integration is $\sim$4x more expensive per swap than CPMM. Acceptable for compute markets (trades are infrequent relative to DeFi), but may not suit high-frequency trading.
    \item \textbf{Parameter sensitivity}: Incorrect $\lambda$ can cause excessive slippage (too high) or instability (too low). Adaptive mechanisms partially address this.
    \item \textbf{Cold start}: New resource types have no price history; initial $\bar{q}_i$ must be set by governance or bootstrapped from external markets.
\end{enumerate}

\subsection{Future Work}

\begin{itemize}[leftmargin=1.1em]
    \item \textbf{Cross-chain liquidity}: Extend HMM to bridge liquidity across multiple chains via IBC or similar protocols.
    \item \textbf{Privacy-preserving execution}: Integrate encrypted TEE attestations to support confidential compute jobs.
    \item \textbf{Quantum extensions}: Replace symplectic integrators with unitary operators for quantum compute resource markets.
    \item \textbf{Reinforcement learning fees}: Train fee parameters via RL to maximize LP returns while maintaining stability.
\end{itemize}

